\begin{blocksection}
\question Write \lstinline{append}, which takes two linked lists and returns a new linked list with all of the elements of the first followed by all of the elements of the second. Then use \lstinline{append} to write \lstinline{reverse}, which returns a new linked list with the elements in reverse order.

\begin{lstlisting}
def append[T](a: LinkedList[T], b: LinkedList[T]) -> LinkedList[T]:
    """
    >>> print(append(Link(1, Link(2)), Link(7, Link(8))))
    (1 2 7 8)
    >>> print(append((), Link(5)))
    (5)
    >>> print(append(Link(1), ()))
    (1)
    >>> append((), ())
    ()
    """
    if ______________________________________:

        return ______________________________

    return ____________________________________________________


def reverse[T](lnk: LinkedList[T]) -> LinkedList[T]:
    """
    >>> print(reverse(Link(1, Link(2, Link(3)))))
    (3 2 1)
    >>> reverse(())
    ()
    """
    if ______________________________________:

        return ______________________________

    return ____________________________________________________
\end{lstlisting}

\begin{solution}[1in]
\begin{lstlisting}
def append[T](a: LinkedList[T], b: LinkedList[T]) -> LinkedList[T]:
    if not isinstance(a, Link):
        return b
    return Link(a.first, append(a.rest, b))

def reverse[T](lnk: LinkedList[T]) -> LinkedList[T]:
    if not isinstance(lnk, Link):
        return ()
    return append(reverse(lnk.rest), Link(lnk.first))
\end{lstlisting}
\end{solution}

\end{blocksection}

\begin{questionmeta}
\textbf{Teaching Tips}\par\nopagebreak[4]
Suggested Time: 10 min; Difficulty: Easy--Medium
\nopagebreak[4]
\begin{itemize}
    \item \textbf{Append:} When \lstinline{a} is empty, the answer is just \lstinline{b}. Otherwise keep \lstinline{a.first} on the front and append \lstinline{b} to the rest of \lstinline{a}.
    \item \textbf{Reverse:} Assume \lstinline{reverse(lnk.rest)} returns the rest reversed. The first element belongs at the end, so it is appended as a one-element list.
    \item \textbf{Types:} \lstinline{Link(lnk.first)} is a one-element linked list, so it is a valid second argument to \lstinline{append}.
\end{itemize}
\end{questionmeta}
